\chapter*{Zusammenfassung}

Diese Ausarbeitung untersucht den Verlauf von Speedrunning-Weltrekorden auf der Plattform \textit{speedrun.com}. Auf Basis von 876~Weltrekord-Datenpunkten aus 17~Spielen in sieben~Genres werden drei Forschungsfragen beantwortet: (1)~Wie stark hat sich die Bestzeit prozentual je Spielkategorie reduziert? (2)~Gibt es einen Sättigungspunkt, ab dem weitere Verbesserungen abnehmen? (3)~Welchen Einfluss haben externe Beschleuniger auf die Rekordentwicklung --- konkret: setzt ein Durchbruch die Abflachungsrate zurück~(3a), ist der Tool-Einfluss im Datensatz messbar~(3b), und wie lange hält ein Weltrekord im Mittel je Genre und Jahrzehnt~(3c)?

Als Methoden kommen nichtparametrische Tests (Kruskal-Wallis, Mann-Whitney~U), der Kaplan-Meier-Schätzer für Überlebenszeitanalysen sowie AIC-basierter Modellvergleich mit anschließendem Chow-Test auf Strukturbrüche zum Einsatz.

Die Ergebnisse zeigen erhebliche Unterschiede in der Zeitreduktion zwischen Spielen (\textbf{9,4~\% bis 99,5~\%}). Verbesserungsdynamiken weisen universelle Strukturbrüche auf (Chow-Test, \textbf{alle 17~Spiele} signifikant), was auf initiale Optimierungsschübe durch frühe Community-Aktivität hindeutet. Nach einem Durchbruch-Ereignis zeigen \textbf{5 von 6} auswertbaren Spielen ein klares Plateau ($\phi < 0{,}5$); Durchbrüche setzen die Abflachungsrate demnach nicht zurück, sondern leiten eine Konsolidierungsphase ein. Für \textbf{2 von 17~Spielen} ist ein stabiler Modell-Grenzwert bestimmbar: \textit{Half-Life~2} hat \textbf{95,5~\%} der möglichen Reduktion bis zum Modell-Floor erreicht, \textit{Pokémon Red/Blue} \textbf{82,5~\%}. Die Lebensdauer von Weltrekorden unterscheidet sich signifikant zwischen Genres ($\mathbf{H(6) = 43{,}82}$; $\mathbf{p < 0{,}001}$; $\eta^2 \approx 0{,}044$), wobei RPG-Rekorde mit einem Kaplan-Meier-Median von \textbf{52~Tagen} deutlich länger Bestand haben als Rekorde anderer Genres (Median: \textbf{7–17~Tage}).

\clearpage
